\SourceLeaf{99}%
einführen wollte, lieber diese als seine Funktion li~$x = \int \frac{dx}{\log x}$ gewählt. Eine
einwertige Funktion ist stets als unvergleichlich klassischer und einfacher anzusehen als
eine mehrwertige, zumal $\log x$ selbst schon eine mehrwertige Funktion ist. $\mathord{.}\,\mathord{.}\,\mathord{.}\,\mathord{.}\,\mathord{.}\,\mathord{.}\,\mathord{.}\,\mathord{.}\,\mathord{.}$

$\mathord{.}\,\mathord{.}\,\mathord{.}\,\mathord{.}$ Im Übrigen glaube ich, dass die Ausdehnung der Untersuchungen auf imaginäre
Argumente zu höchst interessanten Ergebnissen führen wird. Aus den
oben genannten Gründen würde ich allerdings lieber die Funktion $\int \frac{e^x-1}{x}\,dx$ als
$\int \frac{dx}{\log x}$ wählen, weil ich vermute, dass erstere elegantere Ergebnisse liefern
wird. So möchte ich zum Beispiel sehr gern wissen, ob diese Funktion oder, was dasselbe ist,
die Reihe

\[
x + \tfrac14 xx + \tfrac{1}{18}x^3 + \text{etc.}
\]

\noindent für gewisse endliche Werte von $x$ der Form $a+bi$ wohl $0$ werden kann. Mit
Gewissheit kann ich es noch nicht behaupten, obwohl es mir sehr wahrscheinlich ist. Gibt es
solche Werte (dann gewiss unendlich viele), so sind dies sehr bemerkenswerte Größen,
und die ganze Reihe lässt sich in unendlich viele Faktoren der Form $(1+2ax+\beta xx)$ zerlegen.
---

$\mathord{.}\,\mathord{.}\,\mathord{.}\,\mathord{.}\,\mathord{.}\,\mathord{.}\,\mathord{.}\,\mathord{.}\,\mathord{.}\,\mathord{.}\,\mathord{.}\,\mathord{.}\,\mathord{.}$

\Needspace{14\baselineskip}
\begin{center}
BEMERKUNG.
\end{center}

\noindent\hspace*{\parindent}Diese grundsätzliche Erörterung der Theorie der Funktionen einer komplexen Variablen
stammt, wie aus dem Text hervorgeht, aus einem Brief an \textsc{Bessel} vom 18.\ Dezember 1811,
der bereits im »Briefwechsel zwischen \textsc{Gauss} und \textsc{Bessel}«
(Leipzig, 1880) veröffentlicht wurde. Wahrscheinlich war die von
\textsc{Gauss} hier entwickelte Auffassung ihm schon lange vor der Abfassung dieses
Briefes vertraut. Eben aufgrund dieser Auffassung dürfte \textsc{Gauss} zum\,Beispiel\ die
vermutlich einer weit früheren Zeit angehörenden Untersuchungen in der zweiten Hälfte
von Art.\ 12 aus dem »Arithmetisch-geometrischen Mittel« (Gesammelte\ Werke III S.\ 378) angestellt
haben. \hfill \textsc{Fricke}.
